\section{Limitations}\label{sec:limitations}

The FX evidence comes from one market maker, while the weekly reports combine positions across dealers. Related categories and contracts selected for continuous coverage remain dependent aggregate tasks, not independent dealer replications. Reported net-position peaks leave gross portfolios, hedging arrangements and economic capital requirements unmeasured. Downloaded economic dates omit historical first releases, reporting delays and unrecorded excursions between observations.

The original FX evaluation period was inspected before conditional models were specified, and all earlier-period comparisons are retrospective. The comparisons identify predictive performance, leaving causal effects and the contributions of individual features unresolved. Dependence between observations prevents a coverage guarantee for fitted limits. Choosing among the tested penalties after observing evaluation losses provides a diagnostic comparison, not an implementable selection rule. Validation-length and high-quantile checks were designed after the initial results, with numerical failures retained. Only 6 of 10 high-quantile compatibility fits pass, leaving both FX comparisons unavailable under the rejected repair.

The cost comparisons apply specified reserve, breach and exceedance prices to fixed reported limits. Realized trading returns, funding gains, intervention costs and the effect of imposing those limits on subsequent trading remain unmeasured.

\FloatBarrier

\section{Conclusion}\label{sec:conclusion}

Historical rules are hard to beat, and short training or validation histories favor them in the weekly comparisons. Matched FX gains depend on the historical reference, answering RQ1 (Table~\ref{tab:fx}). Longer validation narrows weekly median loss gaps, answering RQ2 (Table~\ref{tab:v13-validation}). The tested selection rule has 0.00\% median weekly excess, yet 8 of 60 cases lose by over ten percent, reaching 80.98\% at worst (Section~\ref{sec:v13-rp}). Long histories support the larger FX gains but do not by themselves secure an advantage (Figure~\ref{fig:v11-history}). For RQ3, higher reserve charges and lower fixed breach charges generally favor regression, so cost structure determines the preferred rule (Figure~\ref{fig:v10-cost}).

Use historical headroom when the available validation history is short, and require convincing comparative evidence before adopting conditional forecasts. Compare methods with the same training observations and update times, and identify shorter rolling references separately. Report forecast scores, breach counts and cost profiles together, since a lower score can favor a limit with more breaches. This joint report lets users select limits according to their own reserve and breach charges.

\FloatBarrier

